Encuentra el valor numérico de Las siguientes expresiones:

        \begin{multicols}{2}
            \begin{parts}
                \part $\large \left( \dfrac{x-y}{a+b} \right)^{3}$ cuando $a=-2$, $b=7$, $x=-6$ y $y=4$.

                \begin{solutionbox}{1.6cm}\footnotesize%
                    $\left( \dfrac{x-y}{a+b} \right)^{3}=\left( \dfrac{-6-4}{-2+7} \right)^{3}=\left( \dfrac{-10}{5} \right)^{3}=\left( -2 \right)^{3}=$ \fillin[-8][0cm]
                \end{solutionbox}

                \part $\large 5m-2n+x$ cuando $m=-3$, $n=4$ y $x=5$.

                \begin{solutionbox}{1.6cm}\footnotesize%
                    $5m-2n+x=5(-3)-2(4)+5=-15-8+5=$ \fillin[-18][0cm]\end{solutionbox}
            \end{parts}
        \end{multicols}
    